% 核心结论：用真实电阻符号画出每条星形臂，中心名移离斜支路
% 草图：[真实几何与节点 | 条件、推导、结论]；17cm印宽，最小字号8pt。
% 图契约：星形中心消元保持任意外端电势下的响应；图中的端口、几何、箭头和方程对应本题。
% 版面：沿用前三模块的几何图+推导分区；目标印宽17cm，简图至中等密度。
% 布局：左侧真实结构，右侧关键关系；真实电路边不作为架构流程箭头处理。
\documentclass[tikz,border=7pt]{standalone}
\usepackage{fontspec,xeCJK,amsmath,amssymb}
\IfFontExistsTF{Times New Roman}{\setmainfont{Times New Roman}}{\setmainfont{TeX Gyre Termes}}
\IfFontExistsTF{Heiti SC}{\setCJKmainfont[AutoFakeBold=1.5]{Heiti SC}}{\setCJKmainfont[AutoFakeBold=1.5]{FandolHei-Regular}}
\usetikzlibrary{arrows.meta,bending,calc,patterns,positioning,decorations.pathmorphing}
\definecolor{ink}{HTML}{243745}
\definecolor{blue}{HTML}{275DAB}
\definecolor{green}{HTML}{177C63}
\definecolor{red}{HTML}{BD4935}
\definecolor{orange}{HTML}{AA6C16}
\definecolor{purple}{HTML}{715393}
\tikzset{every node/.style={font=\fontsize{9.3}{12.5}\selectfont,text=ink},
arrow/.style={-{Stealth[length=5pt 1.5,width'=0pt 3]},shorten >=2pt,shorten <=1pt,line width=1pt},
arrow short/.style={-{Stealth[length=3pt 1.0,width'=0pt 2.5]},shorten >=1pt,shorten <=0pt,line width=.8pt},
axis/.style={arrow,draw=ink!60,line width=.6pt},
vec/.style={arrow,draw=blue},
wire/.style={draw=ink!75,line width=.8pt},
guide/.style={dashed,draw=ink!45,line width=.5pt},
card/.style={rounded corners=2pt,draw=blue!65,fill=blue!4,align=center,text width=6.5cm,inner xsep=9pt,inner ysep=7pt},
title/.style={font=\fontsize{10.5}{14}\selectfont\bfseries,text width=15cm,align=center},
label/.style={font=\fontsize{8.4}{11}\selectfont,align=center},
note/.style={label,text width=6.1cm}}
\begin{document}
\begin{tikzpicture}[x=1cm,y=1cm]
\path[use as bounding box] (-8,-3.5) rectangle(8,3.65);
\node[title] at (0,3.25) {M4-080\quad 星形中心消元保持任意外端电势下的响应};
\begin{scope}[shift={(-4.05,0)}]\draw[wire,draw=ink](0,0)--(4.5924254968025736e-17,0.7499999999999999) (7.960204194457795e-17,1.2999999999999998)--(1.2552629691260368e-16,2.05);\begin{scope}[shift={(6.276314845630184e-17,1.025)},rotate=90.0]\draw[draw=ink,fill=white,line width=.8pt](-0.275,-.085)rectangle(0.275,.085);\end{scope}
\node[label,anchor=east,fill=white,inner sep=1pt]at(-0.32999999999999996,1.025){$r$};\draw[wire,draw=ink](0,0)--(0.6495190528383289,0.3749999999999999) (1.1258330249197703,0.6499999999999998)--(1.7753520777580991,1.0249999999999997);\begin{scope}[shift={(0.8876760388790496,0.5124999999999998)},rotate=29.999999999999996]\draw[draw=ink,fill=white,line width=.8pt](-0.275,-.085)rectangle(0.275,.085);\end{scope}
\node[label,anchor=center,fill=white,inner sep=1pt]at(0.6376760388790497,0.9455127018922191){$r$};\draw[wire,draw=ink](0,0)--(0.6495190528383289,-0.3749999999999999) (1.1258330249197703,-0.6499999999999998)--(1.7753520777580991,-1.0249999999999997);\begin{scope}[shift={(0.8876760388790496,-0.5124999999999998)},rotate=-29.999999999999996]\draw[draw=ink,fill=white,line width=.8pt](-0.275,-.085)rectangle(0.275,.085);\end{scope}
\node[label,anchor=center,fill=white,inner sep=1pt]at(1.1376760388790494,-0.07948729810778049){$r$};\draw[wire,draw=ink](0,0)--(4.5924254968025736e-17,-0.7499999999999999) (7.960204194457795e-17,-1.2999999999999998)--(1.2552629691260368e-16,-2.05);\begin{scope}[shift={(6.276314845630184e-17,-1.025)},rotate=-90.0]\draw[draw=ink,fill=white,line width=.8pt](-0.275,-.085)rectangle(0.275,.085);\end{scope}
\node[label,anchor=west,fill=white,inner sep=1pt]at(0.33000000000000007,-1.025){$r$};\draw[wire,draw=ink](0,0)--(-0.6495190528383289,-0.3749999999999999) (-1.1258330249197703,-0.6499999999999998)--(-1.7753520777580991,-1.0249999999999997);\begin{scope}[shift={(-0.8876760388790496,-0.5124999999999998)},rotate=-150.0]\draw[draw=ink,fill=white,line width=.8pt](-0.275,-.085)rectangle(0.275,.085);\end{scope}
\node[label,anchor=center,fill=white,inner sep=1pt]at(-0.6376760388790497,-0.9455127018922191){$r$};\draw[wire,draw=ink](0,0)--(-0.6495190528383288,0.375) (-1.12583302491977,0.6500000000000001)--(-1.775352077758099,1.0250000000000001);\begin{scope}[shift={(-0.8876760388790494,0.5125000000000001)},rotate=150.0]\draw[draw=ink,fill=white,line width=.8pt](-0.275,-.085)rectangle(0.275,.085);\end{scope}
\node[label,anchor=center,fill=white,inner sep=1pt]at(-1.1376760388790494,0.07948729810778077){$r$};\fill[ink] (1.2552629691260368e-16,2.05)circle(2pt);\node[above=4pt] at(1.2552629691260368e-16,2.05){$1$};
\fill[ink] (1.7753520777580991,1.0249999999999997)circle(2pt);\node[above=4pt] at(1.7753520777580991,1.0249999999999997){$2$};
\fill[ink] (1.7753520777580991,-1.0249999999999997)circle(2pt);\node[below=4pt] at(1.7753520777580991,-1.0249999999999997){$3$};
\fill[ink] (1.2552629691260368e-16,-2.05)circle(2pt);\node[below=4pt] at(1.2552629691260368e-16,-2.05){$4$};
\fill[ink] (-1.7753520777580991,-1.0249999999999997)circle(2pt);\node[below=4pt] at(-1.7753520777580991,-1.0249999999999997){$5$};
\fill[ink] (-1.775352077758099,1.0250000000000001)circle(2pt);\node[above=4pt] at(-1.775352077758099,1.0250000000000001){$6$};
\fill[ink] (0,0)circle(2pt);
\node[label]at(.4,.52){$O$};\end{scope}
\node[card,text width=6.5cm,draw=blue!65,fill=blue!4]at(3.9,1.8){N等臂星形，每臂阻值 $r$};
\node[card,text width=6.5cm,draw=green!65,fill=green!4]at(3.9,0){消元后完全图每边：$R_{ij}^{\text{完}}=Nr$};
\node[card,text width=6.5cm,draw=orange!65,fill=orange!4]at(3.9,-1.8){其他端悬空二端测量：$R_{ij}^{\text{测}}=2r$};

\end{tikzpicture}\end{document}